\Needspace{8\baselineskip}
\Article{272}
Den im vorigen Artikel aufgestellten Satz, dass alle ternären Formen mit gegebener Determinante in eine \emph{endliche} Anzahl von Klassen zerfallen, werden wir mit einer ähnlichen Methode wie bei den binären Formen behandeln. Wir zeigen zunächst, wie jede ternäre Form auf eine einfachere Form reduziert werden kann, und dann, dass die Anzahl der einfachsten Formen (zu denen solche Reduktionen führen) bei jeder gegebenen Determinante endlich ist. Nehmen wir \SourceLeaf{313}allgemein an, es sei eine ternäre Form \(f=\left(\begin{smallmatrix}a,&a',&a''\\b,&b',&b''\end{smallmatrix}\right)\) mit der Determinante \(D\) (verschieden von Null) gegeben, die durch die Substitution \((S)\)
\[
\begin{array}{ccc}
\alpha,&\beta,&\gamma\\
\alpha',&\beta',&\gamma'\\
\alpha'',&\beta'',&\gamma''
\end{array}
\]
in eine äquivalente Form \(g=\left(\begin{smallmatrix}m,&m',&m''\\n,&n',&n''\end{smallmatrix}\right)\) übergeht. Unsere Aufgabe besteht darin, \(\alpha\), \(\beta\), \(\gamma\) usw.\ so zu bestimmen, dass die Form \(g\) einfacher als \(f\) wird. Die den Formen \(f\), \(g\) adjungierten Formen seien jeweils \(\left(\begin{smallmatrix}A,&A',&A''\\B,&B',&B''\end{smallmatrix}\right)\), \(\left(\begin{smallmatrix}M,&M',&M''\\N,&N',&N''\end{smallmatrix}\right)\). und werden mit \(F\), \(G\) bezeichnet. Nach Artikel 269 geht dann \(F\) in \(G\) durch die der Substitution \(S\) adjungierte Substitution über, \(G\) dagegen in \(F\) durch die Substitution, die durch Transposition von \(S\) entsteht. Die Zahl
\[
\alpha\beta'\gamma''+\alpha'\beta''\gamma+\alpha''\beta\gamma'-\alpha''\beta'\gamma-\alpha\beta''\gamma'-\alpha'\beta\gamma''
\]
die entweder \(=+1\) oder \(=-1\) sein muss, werden wir mit \(k\) bezeichnen. Unter diesen Voraussetzungen bemerken wir:

I. Setzt man \(\gamma=0\), \(\gamma'=0\), \(\alpha''=0\), \(\beta''=0\), \(\gamma''=1\), so gilt
\[
\begin{aligned}
m&=a\alpha\alpha+2b''\alpha\alpha'+a'\alpha'\alpha',\\
m'&=a\beta\beta+2b''\beta\beta'+a'\beta'\beta',\\
m''&=a''\\
n&=b\beta'+b'\beta,\\
n'&=b\alpha'+b'\alpha,\\
n''&=a\alpha\beta+b''(\alpha\beta'+\beta\alpha')+a'\alpha'\beta'
\end{aligned}
\]
Außerdem muss \(\alpha\beta'-\beta\alpha'\) entweder \(=+1\) oder \(=-1\) sein. Daraus ist klar, dass die binäre Form \((a,b'',a')\) mit der Determinante \(A''\) durch die Substitution \(\alpha\), \(\beta\), \(\alpha'\), \(\beta'\) in die binäre Form \((m,n'',m')\) mit der Determinante \(M''\) verwandelt wird und ihr daher wegen \(\alpha\beta'-\beta\alpha'=\pm1\) äquivalent ist. Also gilt \(M''=A''\), was sich auch unmittelbar leicht bestätigen lässt. Wenn also \((a,b'',a')\) nicht schon eine einfachste Form ihrer Klasse ist, dürfen wir \(\alpha\), \(\beta\), \(\alpha'\), \(\beta'\) so bestimmen, dass \((m,n'',m')\) eine einfachere Form ist. Aus der Theorie der Äquivalenz binärer Formen folgt dabei leicht, dass dies so geschehen kann, dass \(m\) nicht größer als \(\sqrt{-\frac{4}{3}A''}\) ist, wenn \(A''\) negativ ist, oder nicht größer als \(\sqrt{A''}\), wenn \(A''\) positiv ist, oder \(m=0\), wenn \(A''=0\). Somit lässt sich in allen Fällen der (absolute) Wert von \(m\) gewiss unter oder wenigstens bis auf \(\sqrt{\pm\frac{4}{3}A''}\) herabsetzen. Auf diese Weise wird also die Form \(f\) auf eine andere reduziert, deren erster Koeffizient, wenn möglich, kleiner ist und deren adjungierte Form denselben dritten Koeffizienten hat wie die Form \(F\), die \(f\) adjungiert ist. Darin besteht die \emph{erste Reduktion}.

\SourceLeaf{314}
II. Setzt man dagegen \(\alpha=1\), \(\beta=0\), \(\gamma=0\), \(\alpha'=0\), \(\alpha''=0\), so gilt \(k=\beta'\gamma''-\beta''\gamma'=\pm1\). Die der Substitution \(S\) adjungierte Substitution ist also
\[
\begin{array}{rrr}
\pm1,&0,&0\\
0,&\gamma'',&-\beta''\\
0,&-\gamma',&\beta'
\end{array}
\]
durch die \(F\) in \(G\) übergeht. Somit gilt
\[
\begin{aligned}
m&=a,\quad n'=b'\gamma''+b''\gamma',\quad n''=b'\beta''+b''\beta'\\
m'&=a'\beta'\beta'+2b\beta'\beta''+a''\beta''\beta''\\
m''&=a'\gamma'\gamma'+2b\gamma'\gamma''+a''\gamma''\gamma''\\
n&=a'\beta'\gamma'+b(\beta'\gamma''+\gamma'\beta'')+a''\beta''\gamma''\\
M'&=A'\gamma''\gamma''-2B\gamma'\gamma''+A''\gamma'\gamma'\\
N&=-A'\beta''\gamma''+B(\beta'\gamma''+\gamma'\beta'')-A''\beta'\gamma'\\
M''&=A'\beta''\beta''-2B\beta'\beta''+A''\beta'\beta'
\end{aligned}
\]
Daraus folgt, dass die binäre Form \((A'',B,A')\) mit der Determinante \(Da\) durch die Substitution \(\beta'\), \(-\gamma'\), \(-\beta''\), \(\gamma''\) in die Form \((M'',N,M')\) mit der Determinante \(Dm\) übergeht und ihr daher (wegen \(\beta'\gamma''-\gamma'\beta''=\pm1\) oder wegen \(Da=Dm\)) äquivalent ist. Wenn also \((A'',B,A')\) nicht schon eine einfachste Form ihrer Klasse ist, lassen sich die Koeffizienten \(\beta'\), \(\gamma'\), \(\beta''\), \(\gamma''\) so bestimmen, dass \((M''.\,N,M')\) einfacher ist. Dies ist stets so möglich, dass \(M''\) ohne Rücksicht auf das Vorzeichen nicht größer als \(\sqrt{\pm\frac{4}{3}Da}\) ist. Auf diese Weise wird die Form \(f\) also auf eine andere mit demselben ersten Koeffizienten reduziert, deren adjungierte Form aber, wenn möglich, einen kleineren dritten Koeffizienten hat als die Form \(F\), die \(f\) adjungiert ist. Darin besteht die \emph{zweite Reduktion}.

III. Ist also \(f\) eine ternäre Form, auf die weder die erste noch die zweite Reduktion anwendbar ist, \emph{das heißt}, die durch keine von beiden in eine einfachere Form verwandelt werden kann: so gilt notwendig sowohl \(aa<\) oder \(=\frac{4}{3}A''\) als auch \(A''A''<\) oder \(=\frac{4}{3}aD\), ohne Rücksicht auf das Vorzeichen. Also ist \(a^4\) \(<\) oder \(=\frac{16}{9}A''A''\), somit \(a^4<\) oder \(=\frac{64}{27}aD\), \(a^3<\) oder \(=\frac{64}{27}D\), und \(a<\) oder \(=\frac{4}{3}\sqrt[3]{D}\); daraus wiederum \(A''A''<\) oder \(=\frac{16}{9}\sqrt[3]{D^4}\) und \(A''<\) oder \(=\frac{4}{3}\sqrt[3]{D^2}\). Solange also \(a\) oder \(A''\) diese Grenzen noch überschreitet, muss die eine oder die andere der vorigen Reduktionen auf die Form \(f\) anwendbar sein. --- Dieser Schluss darf jedoch nicht umgekehrt werden, da es durchaus häufig
\SourceLeaf{315}
vorkommt, dass eine ternäre Form, deren erster Koeffizient sowie der dritte Koeffizient ihrer adjungierten Form schon unter jenen Grenzen liegen, dennoch durch die eine oder die andere Reduktion weiter vereinfacht werden kann.

IV. Wendet man aber auf eine beliebige gegebene ternäre Form mit der Determinante \(D\) abwechselnd die erste und die zweite Reduktion an, \emph{das heißt}, auf sie selbst die erste oder zweite, auf die daraus entstandene Form die zweite oder erste, auf die nächste daraus entstandene Form wieder die erste oder zweite usw., so muss man offenbar schließlich zu einer Form gelangen, auf die keine der beiden mehr anwendbar ist. Denn sowohl die absoluten Beträge der ersten Koeffizienten der so entstehenden Formen als auch die der dritten Koeffizienten ihrer adjungierten Formen bleiben fortwährend abwechselnd gleich und nehmen ab. Dieser Vorgang muss deshalb schließlich irgendwo enden, da man sonst zwei unendliche Folgen fortwährend abnehmender Zahlen hätte. Damit haben wir einen bemerkenswerten Satz gewonnen: \emph{Jede ternäre Form mit der Determinante \(D\) kann auf eine andere äquivalente Form reduziert werden, deren erster Koeffizient nicht größer als \(\frac{4}{3}\sqrt[3]{D}\) und der dritte Koeffizient ihrer adjungierten Form nicht größer als \(\frac{4}{3}\sqrt[3]{D^2}\) ist, ohne Rücksicht auf das Vorzeichen, sofern die gegebene Form diese Eigenschaften nicht schon selbst besitzt.} --- Statt des ersten Koeffizienten der Form \(f\) und des dritten der zu \(f\) adjungierten Form hätten wir auf genau dieselbe Weise entweder den ersten Koeffizienten der Form selbst und den zweiten ihrer adjungierten Form behandeln können; oder den zweiten der Form selbst und den ersten oder dritten ihrer adjungierten Form; oder den dritten der Form selbst und den ersten oder zweiten ihrer adjungierten Form. Auf allen diesen Wegen wären wir ebenso zum vorgesehenen Ziel gelangt. Es ist aber zweckmäßig, beständig bei einer Methode zu bleiben, damit die zugehörigen Operationen leichter in einen festen Algorithmus gefasst werden können. Schließlich bemerken wir, dass sich für die beiden Koeffizienten, die wir unter feste Grenzen herabzusetzen gelehrt haben. noch kleinere Grenzen bestimmen lassen, wenn man die definiten von den indefiniten Formen trennt; für unser gegenwärtiges Vorhaben ist dies jedoch nicht nötig.
